\chapter{Lesson 7: Visiting Friends}

\section{Vocabulary}

\begin{center}
\begin{tabular}{llll}
\toprule
Characters & Pinyin & Part of speech & Meaning \\
\midrule
\zh{进} & jìn & verb & to enter \\
\zh{快} & kuài & adverb / adjective & quickly; fast \\
\zh{进来} & jìnlai & verb + complement & to come in \\
\zh{介绍} & jièshào & verb & to introduce \\
\zh{一下} & yíxià & measure phrase & (softens a verb) once, a bit \\
\zh{高兴} & gāoxìng & adjective & happy, glad \\
\zh{漂亮} & piàoliang & adjective & pretty \\
\zh{坐} & zuò & verb & to sit \\
\zh{在} & zài & verb / preposition & to be at; at, in \\
\zh{哪儿} & nǎr & question word & where \\
\zh{学校} & xuéxiào & noun & school \\
\zh{喝} & hē & verb & to drink \\
\zh{茶} & chá & noun & tea \\
\zh{咖啡} & kāfēi & noun & coffee \\
\zh{水} & shuǐ & noun & water \\
\zh{可乐} & kělè & noun & cola \\
\zh{啤酒} & píjiǔ & noun & beer \\
\zh{吧} & ba & particle & (suggestion / supposition) \\
\zh{要} & yào & verb & to want \\
\zh{瓶} & píng & measure word & bottle \\
\zh{杯} & bēi & measure word & cup, glass \\
\zh{给} & gěi & verb & to give \\
\zh{玩儿} & wánr & verb & to have fun, to hang out \\
\zh{图书馆} & túshūguǎn & noun & library \\
\zh{聊天(儿)} & liáo tiān(r) & verb-object & to chat \\
\zh{才} & cái & adverb & not until, only then \\
\zh{回家} & huí jiā & verb-object & to go home \\
\zh{了} & le & particle & (completed action) \\
\bottomrule
\end{tabular}
\end{center}

\section{Etiquette of a visit}

\begin{aside}
When visiting a Chinese home, bring a small gift (fruit, tea, something from your home
town). The host will say \zh{你太客气了！} ``You're too polite!'' before accepting, and
may set it aside without opening it. Guests are offered tea or a drink straight away;
refusing once out of politeness is normal, and a good host offers again. Shoes usually
come off at the door. When leaving, the host walks the guest out and the guest says
\zh{请留步} \emph{qǐng liú bù} ``please don't trouble to see me out''.
\end{aside}

\section{Grammar}

\subsection{\zh{一下} and reduplicated verbs}

\zh{一下} after a verb makes a request sound casual and brief, like ``take a quick
look''.
\begin{itemize}
  \item \zh{我介绍一下，这是我的同学小李。} Let me introduce: this is my classmate Xiao Li.
  \item \zh{你看一下这张照片。} Have a look at this photo.
\end{itemize}
Reduplicating a one-syllable verb has the same effect: \zh{看看}, \zh{坐坐},
\zh{聊聊天}.

\subsection{\zh{在}: verb and preposition}

As a verb, \zh{在} means ``to be (located) at''. As a preposition, \zh{在} + place
comes \emph{before} the verb to say where an action happens.

\begin{center}
\begin{tikzpicture}[w/.style={draw, rounded corners, minimum height=0.8cm, align=center, font=\small}]
  \node[w, fill=blue!8] (s) at (0,0) {\zh{我姐姐}\\ subject};
  \node[w, fill=green!12] (z) at (2.4,0) {\zh{在学校}\\ \zh{在} + place};
  \node[w, fill=orange!15] (v) at (4.8,0) {\zh{工作}\\ verb};
  \draw[->, thick] (s) -- (z);
  \draw[->, thick] (z) -- (v);
  \node[font=\small, anchor=west] at (5.8,0) {My older sister works at a school.};
  \node[w, fill=blue!8] (s2) at (0,-1.6) {\zh{他}\\ subject};
  \node[w, fill=green!12] (z2) at (2.4,-1.6) {\zh{在}\\ verb};
  \node[w, fill=orange!15] (v2) at (4.8,-1.6) {\zh{图书馆}\\ place};
  \draw[->, thick] (s2) -- (z2);
  \draw[->, thick] (z2) -- (v2);
  \node[font=\small, anchor=west] at (5.8,-1.6) {He is in the library.};
\end{tikzpicture}
\end{center}

\begin{important}[Place before verb]
English puts ``where'' after the verb (``work at a school''); Chinese puts the
\zh{在}-phrase before it. *\zh{我工作在学校} is a typical error.
\end{important}

Questions: \zh{你在哪儿工作？} ``Where do you work?''; \zh{你家在哪儿？} ``Where is your
home?''

\subsection{\zh{吧}: suggestions and assumptions}

\begin{itemize}
  \item Suggestion: \zh{我们喝茶吧。} ``Let's have tea.''
  \item Softened request: \zh{你坐吧。} ``Have a seat.''
  \item Checking an assumption: \zh{你是小王的姐姐吧？} ``You must be Xiao Wang's older
        sister, right?''
\end{itemize}

\subsection{Offering a drink: \zh{要}, \zh{还是}, \zh{给}}

\begin{itemize}
  \item \zh{你要喝什么？要茶还是要咖啡？} What would you like to drink, tea or coffee?
  \item \zh{给我一杯水吧。} Give me a glass of water, please.
  \item \zh{对不起，我们没有咖啡了。} Sorry, we're out of coffee.
\end{itemize}

\subsection{The particle \zh{了} (completed action)}

\zh{了} after a verb indicates that an action has been completed or has taken place.

\begin{itemize}
  \item \zh{昨天我喝了两杯茶。} Yesterday I drank two cups of tea.
  \item \zh{她买了一本书。} She bought a book.
  \item \zh{我们昨天晚上在小李家聊天，聊到十一点。}
\end{itemize}

\begin{important}[\zh{了}: first rules]
\begin{itemize}
  \item With an object, \zh{了} after the verb normally needs a quantified or specific
        object (\zh{喝了两杯茶}); a bare object (\zh{我喝了茶}) sounds unfinished.
        Alternatively put \zh{了} at the end: \zh{我喝茶了}.
  \item The negative uses \zh{没(有)} and \emph{drops} \zh{了}: \zh{我没喝茶}, not
        *\zh{我没喝了茶}.
  \item Questions: \zh{你喝茶了吗？} or \zh{你喝没喝茶？}/\zh{你喝茶了没有？}
  \item \zh{了} is \emph{not} a past tense. Habitual past actions take no \zh{了}:
        \zh{以前我常常去图书馆}.
\end{itemize}
\end{important}

\subsection{\zh{才} ``not until''}

\zh{才} before a verb says the action happened later than expected.
\begin{itemize}
  \item \zh{我们聊天聊到十二点，我才回家。} We chatted until twelve; only then did I go
        home.
  \item \zh{他八点上课，八点十分才来。} His class is at eight; he didn't come until
        8:10.
\end{itemize}
Note: sentences with \zh{才} do \emph{not} take \zh{了}. (Contrast with \zh{就} in Lesson
10.)

\section{Model conversation}

\begin{itemize}
  \item \textbf{A:} \zh{请进，请进！快进来！}
  \item \textbf{B:} \zh{这是我给你买的水果。}
  \item \textbf{A:} \zh{你太客气了。来，我介绍一下：这是我的室友，王红。}
  \item \textbf{B:} \zh{你好，认识你很高兴。}
  \item \textbf{C:} \zh{认识你我也很高兴。请坐。你在哪儿工作？}
  \item \textbf{B:} \zh{我不工作，我在滑铁卢大学学中文。}
  \item \textbf{A:} \zh{你们喝点儿什么？茶还是可乐？}
  \item \textbf{B:} \zh{我喝茶吧。谢谢。}
\end{itemize}
(\zh{室友} shìyǒu = roommate; \zh{水果} shuǐguǒ = fruit; \zh{认识你很高兴} = nice to
meet you.)

\section*{Review}
\begin{itemize}
  \item \zh{一下} or a reduplicated verb softens a request.
  \item \zh{在} + place goes \emph{before} the verb.
  \item \zh{吧} for suggestions and ``I suppose''.
  \item V + \zh{了} marks a completed action; negate with \zh{没} and no \zh{了}.
  \item \zh{才} = later than expected, no \zh{了}.
  \item Visiting: bring a gift, expect to be offered tea, and expect polite refusals.
\end{itemize}
